\section{Závěr}

Obrázky nesou o popularitě hry informaci, která v metadatech není. Oba hybridní modely mají
na validační i testovací množině nižší MAE než všechny baseline nad metadaty a ablace ukázala,
že kombinace obou modalit je lepší než každá zvlášť. Největší přínos měla volba předtrénovaného
encoderu: CLIP snížil testovací MAE modelu B z 1,183 na 1,108. Rozdíl mezi vlastní CNN
(model A) a zmrazeným ConvNeXt-Tiny (model B) je malý a při třech seedech jen orientační.

Encodery modelu B byly předtrénovány na velkém množství dalších dat s mnohem větším výpočetním
rozpočtem, nejde tedy o srovnání se stejným rozpočtem jako u modelu A. Hlavní omezení studie
jsou hledání hyperparametrů s jedním seedem a jen pro encoder ConvNeXt, limit 15 epoch, který
pravděpodobně omezil model A, a výstup modelů v bfloat16, který predikce zaokrouhluje
přibližně na 0,03. Všechny modely stahují predikce ke středu rozdělení a velmi populární hry
systematicky podhodnocují. Publikované výsledky nejsou přímo srovnatelné, protože se liší
čištění dat, definice cílové proměnné i rozdělení dat.
% TODO: ověřit, zda existují publikované výsledky na tomto datasetu

\subsection*{Využití generativní AI}

% TODO:
